\documentclass[12pt,a4paper]{article}

% Encoding and fonts — XeLaTeX/LuaLaTeX
\usepackage{fontspec}
\newfontfamily\hebrewfont[Script=Hebrew]{Noto Serif Hebrew}
\newcommand{\heb}[1]{{\hebrewfont #1}}
\newfontfamily\igfont[Ligatures=TeX]{Noto Serif}
\newcommand{\igtext}[1]{{\igfont #1}}

% Page layout
\usepackage[top=1in, bottom=1in, left=1in, right=1in]{geometry}

% Typography
\usepackage{microtype}
\usepackage{parskip}

% Language
\usepackage[english]{babel}

% Headers and footers
\usepackage{fancyhdr}
\pagestyle{fancy}
\fancyhf{}
\fancyhead[L]{\leftmark}
\fancyhead[R]{\thepage}
\fancyfoot[C]{\thepage}
\renewcommand{\headrulewidth}{0.4pt}
\renewcommand{\footrulewidth}{0pt}

% Section styling
\usepackage{titlesec}
\titleformat{\section}{\Large\bfseries}{\thesection}{1em}{}
\titleformat{\subsection}{\large\bfseries}{\thesubsection}{1em}{}
\titleformat{\subsubsection}{\normalsize\bfseries}{\thesubsubsection}{1em}{}
\titlespacing*{\section}{0pt}{2.5ex plus 1ex minus .2ex}{1.5ex plus .2ex}
\titlespacing*{\subsection}{0pt}{2ex plus 1ex minus .2ex}{1ex plus .2ex}
\titlespacing*{\subsubsection}{0pt}{1.5ex plus 1ex minus .2ex}{0.8ex plus .2ex}

% List formatting
\usepackage{enumitem}
\setlist[itemize]{topsep=0.5ex, itemsep=0.25ex, parsep=0pt}
\setlist[enumerate]{topsep=0.5ex, itemsep=0.25ex, parsep=0pt}

% Essential packages
\usepackage{graphicx}
\usepackage[table]{xcolor}
\usepackage{booktabs}
\usepackage{array}
\usepackage{tabularx}
\usepackage{longtable}
\usepackage{amsmath}
\usepackage{amssymb}
\usepackage[normalem]{ulem}
\rowcolors{2}{gray!10}{white}
\renewcommand{\arraystretch}{1.3}

% Cross-references
\usepackage{hyperref}
\usepackage{cleveref}

% Custom commands
\newcommand{\pilcrow}{\S}

% Hyperref setup
\hypersetup{
    colorlinks=true,
    linkcolor=blue,
    filecolor=magenta,
    urlcolor=cyan,
}

% Code listings
\usepackage{listings}
\definecolor{codebg}{RGB}{248,248,248}
\definecolor{codeframe}{RGB}{200,200,200}
\definecolor{codekeyword}{RGB}{0,0,180}
\definecolor{codecomment}{RGB}{0,128,0}
\definecolor{codestring}{RGB}{163,21,21}
\lstset{
    basicstyle=\ttfamily\small,
    breaklines=true,
    breakatwhitespace=false,
    frame=single,
    framerule=0.5pt,
    rulecolor=\color{codeframe},
    backgroundcolor=\color{codebg},
    keepspaces=true,
    numbers=left,
    numberstyle=\tiny\color{gray},
    numbersep=8pt,
    showstringspaces=false,
    tabsize=4,
    xleftmargin=1em,
    xrightmargin=1em,
    aboveskip=1em,
    belowskip=1em,
    keywordstyle=\color{codekeyword}\bfseries,
    commentstyle=\color{codecomment}\itshape,
    stringstyle=\color{codestring},
    literate={_}{{\textunderscore}}1
             {-}{{\textminus}}1
             {<}{{\textless}}1
             {>}{{\textgreater}}1
             {~}{{\textasciitilde}}1
             {^}{{\textasciicircum}}1
}

% YAML language definition for listings
\lstdefinelanguage{YAML}{
    keywords={true,false,null,y,n},
    keywordstyle=\color{blue}\bfseries,
    sensitive=false,
    comment=[l]{\#},
    morecomment=[s]{/*}{*/},
    commentstyle=\color{gray}\ttfamily,
    stringstyle=\color{red}\ttfamily,
    morestring=[b]'',
    morestring=[b]'',
}

% TOML language definition for listings
\lstdefinelanguage{TOML}{
    keywords={true,false},
    keywordstyle=\color{blue}\bfseries,
    sensitive=true,
    comment=[l]{\#},
    commentstyle=\color{gray}\ttfamily,
    stringstyle=\color{red}\ttfamily,
    morestring=[b]'',
    morestring=[b]'',
}

% Dockerfile language definition for listings
\lstdefinelanguage{Dockerfile}{
    keywords={FROM,RUN,CMD,LABEL,EXPOSE,ENV,ADD,COPY,ENTRYPOINT,VOLUME,USER,WORKDIR,ARG,ONBUILD,STOPSIGNAL,HEALTHCHECK,SHELL},
    keywordstyle=\color{blue}\bfseries,
    sensitive=false,
    comment=[l]{\#},
    commentstyle=\color{gray}\ttfamily,
    stringstyle=\color{red}\ttfamily,
    morestring=[b]'',
    morestring=[b]'',
}

% JSON language definition for listings
\lstdefinelanguage{JSON}{
    keywords={true,false,null},
    keywordstyle=\color{blue}\bfseries,
    sensitive=true,
    commentstyle=\color{gray}\ttfamily,
    stringstyle=\color{red}\ttfamily,
    morestring=[b]'',
}

% Code listing float environment
\usepackage{float}
\newfloat{listing}{htbp}{lol}[section]
\floatname{listing}{Listing}

% Admonition boxes
\usepackage{tcolorbox}
\tcbuselibrary{skins,breakable}

% Admonition style definitions
\newtcolorbox{admonitionbox}[2][]{
    enhanced,
    breakable,
    colback=#2!5,
    colframe=#2!80!black,
    fonttitle=\bfseries,
    title=#1,
    left=2mm,
    right=2mm,
}

% Synthon tuple boxes
\newtcolorbox{synthonbox}{
    enhanced,
    colback=blue!5,
    colframe=blue!75!black,
    fontupper=\large,
    center upper,
    boxrule=1pt,
    arc=4pt,
}

\begin{document}

\title{From Adolescence to Adulthood}
\date{\today}

\maketitle

\textbf{Purpose} --- This assignment traces how experiences during adolescence shape identity development, decision-making, and future goals in early adulthood. The analysis draws on Sociocultural Theory, Behaviorism, Social Learning Theory, Operant Conditioning, and Cognitive Theories of Multiple Intelligences and Information Processing.

\begin{center}
\rule{0.5\textwidth}{0.4pt}
\end{center}

I graduated high school in 2022. I write this now --- not then --- because the distance between those two moments is, itself, the subject of the question.

My stated goal is to become an elementary school teacher and work with young children. I will try to explain how I arrived here. The honest answer is: not as linearly as the prompts below might suggest. Early drafts of this reflection listed Bandura and moved on. A second draft felt wrong because it made the path sound inevitable. It wasn''t. It remains, in some ways, unfinished --- which is precisely what James Marcia''s model would predict, and which is what I need to sit with before answering the questions in order.

\begin{center}
\rule{0.5\textwidth}{0.4pt}
\end{center}

\subsection{I. What I Saw Before I Knew What I Was Looking For}

Bandura''s Social Learning Theory states that people learn through observation and imitation. The standard application would be: I watched good teachers, I wanted to be one. That is what I wrote first. It is not wrong --- but it is incomplete, and treating it as complete does a disservice to both the theory and the experience.

The teachers I remember --- really remember, not just recall for an assignment --- were not uniformly patient. Some days were difficult. What I observed was not perfection but \emph{attempts}: adults choosing, repeatedly, to re-explain a concept to a child who was lost, to redirect a classroom that had tipped into noise, to stay in the room when leaving would have been easier. What Bandura captures is not the imitation of a role model''s success but the internalization of their \emph{persistence}. I did not decide to teach because I saw someone teach well. I decided because I saw someone teach \emph{when they didn''t have to}, and something in me registered: that is a life worth choosing.

A counterargument worth entertaining: maybe I am retroactively organizing scattered memories into a coherent origin story. This is the reconstruction bias, well-documented in developmental psychology. I cannot fully discount it. But even if the narrative is partly constructed, the construction itself is data --- the mind does not invent origins from nothing. It selects from what was there.

\begin{center}
\rule{0.5\textwidth}{0.4pt}
\end{center}

\subsection{II. Still Choosing --- Marcia''s Moratorium, Lived Rather Than Categorized}

James Marcia''s identity status model names four positions: diffusion (no exploration, no commitment), foreclosure (commitment without exploration), moratorium (active exploration, no firm commitment), and achievement (commitment following exploration). I would describe myself as in moratorium --- and I say this with the caveat that the label fits \emph{because} I haven''t fully committed, not because I''ve explored exhaustively and landed provisionally. The distinction matters.

Three experiences shaped this state:

\begin{enumerate}
    \item \textbf{Helping classmates.} There were moments in high school --- not planned, not part of any peer-tutoring program --- when I found myself explaining an assignment to someone who was stuck. The satisfaction I felt was not vanity. It was the quiet recognition that I could \emph{make something legible} for another person. I did not immediately think ''''teaching.'''' The thought came later, as a pattern emerged.
    \item \textbf{Watching teachers, again.} The same teachers whose persistence I described above also did something I had not noticed at the time: they calibrated. They read a room. They adjusted. What I am only now articulating is that this calibration is itself a skill --- cognitive, social-emotional, even physical --- and wanting to develop it is different from wanting to \emph{be} those teachers. The model is not a template; it is a provocation.
    \item \textbf{Being around younger children.} This is the hardest to pin down because it is not a single event. It is the accumulation of occasions --- family gatherings, neighborhood kids, moments where I found myself drawn into conversation with someone half my age because they were interesting, not because I was being dutiful. The pull was real. But wanting to be \emph{around} children and wanting to \emph{teach} them are related and not identical. I am still untangling this.
\end{enumerate}

\begin{center}
\rule{0.5\textwidth}{0.4pt}
\end{center}

\subsection{III. What I Bring to the Classroom --- Developmental Domains Reconsidered}

The prompt asks me to map my characteristics onto three domains. I will do so, but I want to resist the temptation to list strengths as though they were completed. They are works in progress --- which is, I think, the more honest answer.

\textbf{Physical.} I have the energy to move through a full school day with young children. This is not a trivial claim: elementary classrooms are physically demanding spaces. But more precisely, what I have is \emph{endurance} --- not the kind that comes from fitness training, but the kind that comes from caring about something enough to stay standing. Stamina, not speed. I can handle noise, motion, the constant low-grade chaos of a room full of bodies who are still learning how to sit still. Whether that stamina will hold up to the reality of 25 students rather than the 3 or 4 I''m used to at family gatherings --- that remains to be seen, honestly.

\textbf{Cognitive.} I can break down complex ideas into simpler ones. But I want to be careful: being able to explain something to \emph{one} confused classmate is not the same as being able to differentiate instruction for thirty students with thirty different cognitive profiles. What I actually have is a starting point --- a facility for translation between ways of thinking --- and a willingness to keep translating until something connects. Gardner''s Multiple Intelligences theory is worth invoking here, though not naively: the categories are useful heuristics, not fixed boxes.

\textbf{Social-Emotional.} Patience and empathy are often listed as ''''soft skills'''' and treated as pleasant additions. In a classroom, they are structural. A child who has been told they are slow to learn will not hear your lesson until they first believe you do not think they are slow. That belief is conveyed through micro-interactions --- tone, waiting time, who you call on, what you do with a wrong answer. I am getting better at reading these moments in real time. I am not yet good at them. The gap between ''''getting better'''' and ''''good'''' is where my career will happen, and I need to be comfortable living in that gap for a while.

\begin{center}
\rule{0.5\textwidth}{0.4pt}
\end{center}

\subsection{IV. On Paths Not Taken --- Operant Conditioning}

The prompt asks whether I have tried other majors or work experiences, and to use Operant Conditioning. I have not tried other majors. One interpretation: this is foreclosed commitment. Marcia would say no. But Skinner would say something worth hearing: reinforcement is not only about what you \emph{have} experienced. It is also about what you have been \emph{rewarded for}.

My entire academic trajectory has been shaped by positive reinforcement: praise for helping others, encouragement toward service-oriented roles, a social environment that validates caring professions. I have not explored alternatives --- but perhaps more precisely, I have not been \emph{punished} for this path nor \emph{rewarded} for alternatives, so the reinforcement schedule has, by default, narrowed my field of view. If I had been praised equally for analytical or technical work, I might be in a different major. This is not an accusation and not regret. It is simply the behavioral fact: the absence of competing reinforcement is itself a form of shaping.

\begin{center}
\rule{0.5\textwidth}{0.4pt}
\end{center}

\subsection{V. If the Path Closed --- What Then?}

If I could not teach elementary school, I would look toward child development or counseling. Two criteria would guide that pivot:

\textbf{First: transferable strengths.} What can I actually \emph{do}? Patience. Communication. The ability to sit with someone who is struggling and not rush to fix them. These are not uniquely teaching skills. They are human skills, and careers built on human skills tend to be adjacent. Child development and counseling fit. So would social work. So would a dozen other things. The exercise of naming a ''''next choice'''' reveals something uncomfortable: my identity is more diffuse than I admitted in Section II. The teaching goal is a focal point, but it sits on a broader foundation. If the focal point shifted, the foundation would hold --- and I should acknowledge that.

\subsection{\textbf{Second: genuine interest.} Competence without interest is a slow way to lose your life. This is not a theoretical concern. The helping professions are saturated with people who are good at their jobs and quietly miserable because they mistook approval for purpose. I would need to distinguish between ''''I am good at this and people tell me so'''' and ''''this is work I would do even if no one told me I was good at it.'''' Only the second survives the realities of bureaucracy, limited resources, and the emotional toll of caring for children whose circumstances you cannot fix.}

\subsection{VI. Three Goals --- Not a Wish List, a Trajectory}

When asked to picture an ideal future, I listed three goals. I want to present them not as independent items but as layers of the same structure.

\textbf{The work itself.} I want to teach. Specifically, I want to build a classroom where students feel safe enough to be wrong. The connection between error and learning is not automatic for children; it has to be cultivated. Creating that culture is harder than delivering curriculum. It requires a teacher who has made peace with their own mistakes. I am working on that part.

\textbf{Stability.} ''''Financially independent, steady job, secure day-to-day life.'''' These sound mundane. But for a helping career, stability is not separate from the work --- it is the condition that makes the work sustainable. Burnout in teaching is not primarily about the children; it is about resource scarcity, institutional instability, and the emotional burden of caring in a system that does not care back. I want stability not as an escape from the work but as armor that lets me stay in it.

\textbf{Continuous growth.} ''''I want to keep growing as a person.'''' I wrote this originally as a generic aspiration. But in the context of everything above, it takes a specific shape: I want the version of me at thirty to look back at the version of me now not with embarrassment but with gratitude that I was willing to not know. Professional growth and personal growth merge here. You cannot become a better teacher without becoming a more honest human being. The two are structurally identical.

\begin{center}
\rule{0.5\textwidth}{0.4pt}
\end{center}

\subsection{VII. Gardner''s Intelligences --- Applied with the Same Honesty}

Howard Gardner identified multiple intelligences. The prompt asks me to map mine. I will do so, but with the same caveat I offered earlier: these categories describe \emph{tendencies}, not capacities.

\textbf{Interpersonal intelligence.} This is the ability to read others --- their moods, their needs, their unspoken questions. It maps directly onto teaching, but it also maps onto life. I have it, or I am developing it. The distinction matters: if I treat it as a fixed asset, I will stop sharpening it. If I treat it as a practice, it will keep growing. Children detect the difference instantly.

\textbf{Linguistic intelligence.} The ability to use language effectively --- to explain, to frame, to translate. This is what I do when I help a classmate. It is what I am doing in this document. But linguistic intelligence in the classroom is not eloquence. It is \emph{clarity at the level of the listener}. That is a harder standard than it sounds, and it requires surrendering your own vocabulary to find the one that works for the student.

\subsection{\textbf{Intrapersonal intelligence.} Self-knowledge. This is the most important of the three for me, and the one I trust least. I can name my strengths --- patience, communication, empathy. But naming them is not the same as inhabiting them honestly. Intrapersonal intelligence, at its best, means knowing what you are avoiding. Right now, what I might be avoiding is the gap between the teacher I imagine and the teacher I am. Acknowledging that gap is, or should be, the first act of genuine self-awareness.}

\subsection{VIII. Scaffolding --- From Whom, and Toward What}

Vygotsky''s concept of scaffolding describes support that is gradually removed as competence increases. Three forms of scaffolding would serve me:

\textbf{Mentorship from experienced educators.} I need someone who has been in a classroom longer than I have been alive --- or close to it --- who can tell me what I am about to get wrong that no textbook will prepare me for. This scaffolding comes from professors, yes, but more importantly from practicing teachers who can say: the lesson plan will fail, and that is not your failure; it is the lesson plan''s. The mentor''s role is not to give me confidence but to reframe what confidence means when the room is falling apart.

\textbf{Structured field experience.} Student teaching is the obvious scaffold here, but I want to be specific: I need to be in a classroom I did \emph{not} choose, with a mentor teacher whose style I am not guaranteed to like. If student teaching only confirms what I already believe, it is not scaffolding --- it is rehearsal. The value is in being placed somewhere unfamiliar and learning how to adapt. This requires trust from the placement coordinators and humility from me.

\textbf{Emotional support from my community.} Family and friends who encourage the work --- not just the idea of the work. There is a difference between ''''I''m proud of you for wanting to help children'''' and ''''I can see how hard today was, and I am here.'''' The first feels good. The second makes it possible to return tomorrow. This support comes from those closest to me, and its value accumulates quietly over years. I do not want to take it for granted.

\begin{center}
\rule{0.5\textwidth}{0.4pt}
\end{center}

\emph{What this exercise has made clear --- and what I did not know at the start --- is that the question ''''What do you want to be?'''' is not the same as ''''Who are you becoming?'''' The first expects an answer. The second expects a practice. I am practicing.}

\begin{center}
\rule{0.5\textwidth}{0.4pt}
\end{center}

\emph{Structural type: $\langle D_{\text{invomega}};\ T_{\text{bullseye}};\ R_{\text{lyoghlig}};\ P_{\text{pipevar}};\ F_{\text{hardsign}};\ K_{\text{schwa}};\ G_{\text{revapostrophe}};\ \Gamma_{\text{secstress}};\ \Phi_{\text{ctyogh}};\ H_2;\ n{:}m;\ \Omega_{\text{crtwo}} \rangle$}

\emph{Eight promotions applied relative to \texttt{adolescence\_draft} baseline ($\langle D_{\text{invomega}};\ T_{\text{nrleg}};\ R_{\text{lyoghlig}};\ P_{\text{aolig}};\ F_{\text{beltl}};\ K_{\text{turnm}};\ G_\text{gimel};\ \Gamma_{\text{corner}};\ \Phi_{\text{ctyogh}};\ H_0;\ n{:}m;\ \Omega_{\text{closeepsilon}} \rangle$): H₀$\rightarrow$H₂ via narrative of wrong-turn-before-right-turn; $\Gamma$$\land$$\rightarrow$$\Gamma$\_seq via section-to-section necessity; T\_net$\rightarrow$T\_bowtie via authorial surprise and the object speaking back; P\_asym$\rightarrow$P\_pm via named counterarguments and uncertainty tracking; $F_{\text{beltl}}$$\rightarrow$F\_ħ via demonstration over restatement; K\_mod$\rightarrow$K\_slow via unresolved hardest claims; G\_gimel$\rightarrow$G\_ℵ via genuine open questions at close; $\Omega$₀$\rightarrow$$\Omega$\_ℤ₂ via closing echo of the opening at higher resolution.}


\end{document}
